% gabarit-rapport.tex - Gabarit de rapport professionnel (Scriptorium)
%
% Portions adaptees du projet openscience (Synthetic Sciences, InkVell Inc.), Apache-2.0,
% github.com/synthetic-sciences/openscience (scientific_report.sty, scientific_report_template.tex,
% skill scientific-writing). Modifications Marius Yvard, MIT.
%
% Simplifie par rapport a la source : un seul fichier autonome (pas de .sty separe), cinq
% encadres semantiques au lieu de neuf, trois macros statistiques au lieu d'une quinzaine.
% Voir livrer/references/document.md ("Sortie LaTeX") pour la voie complete et le choix entre
% LaTeX, Word, PDF et HTML.
%
% Compilation : xelatex gabarit-rapport.tex (deux passes pour la table des matieres,
% la table des figures, la liste des tableaux et les renvois \ref).
% Avec bibliographie : xelatex, bibtex, xelatex, xelatex (fournir un fichier references.bib,
% voir scripts/citations.py pour le produire depuis des sources verifiees).
%
% Les couleurs et polices du bloc CHARTE ci-dessous sont un jeu de valeurs sobres par defaut.
% Pour appliquer une charte graphique reelle, remplacer tout le bloc entre les marqueurs par
% la sortie de :
%   python3 scripts/theme.py charte-graphique.json --format latex
% (meme logique que --format css pour la sortie HTML : un bloc genere a coller tel quel).
%
\documentclass[11pt,a4paper]{article}

% --- Paquets ---
\usepackage{fontspec}
\usepackage[a4paper,margin=2.5cm,headheight=26pt]{geometry}
\usepackage{xcolor}
\usepackage[most]{tcolorbox}
\usepackage{fancyhdr}
\usepackage{titlesec}
\usepackage{booktabs}
\usepackage{amsmath}
\usepackage{graphicx}
\usepackage{float}
\usepackage[colorlinks=true]{hyperref}
\usepackage[numbers,sort&compress]{natbib}
% caption apres hyperref : l'ordre inverse casse les ancres de legende.
% subcaption charge caption et sert aux figures composees (a) (b).
\usepackage{caption}
\usepackage{subcaption}

% Dossiers ou chercher les images, en plus du dossier du .tex. Les figures SVG
% produites par scripts/figures.py se convertissent en PDF avant inclusion
% (rsvg-convert, inkscape ou cairosvg) : xelatex n'insere pas de SVG.
\graphicspath{{figures/}{images/}}

% ============================================================================
% >>> DEBUT BLOC CHARTE (remplacer par : python3 scripts/theme.py CHARTE.json --format latex)
% ============================================================================
\definecolor{ScriptoriumEncre}{HTML}{2E2A26}
\definecolor{ScriptoriumTrait}{HTML}{8A8175}
\definecolor{ScriptoriumFond}{HTML}{FFFFFF}
\definecolor{ScriptoriumAccent}{HTML}{6E6356}
\definecolor{ScriptoriumPalUn}{HTML}{F4F1EC}
\definecolor{ScriptoriumPalDeux}{HTML}{EEF2F4}
\definecolor{ScriptoriumPalTrois}{HTML}{F1F0EA}
\definecolor{ScriptoriumPalQuatre}{HTML}{F4EEF0}
\setmainfont{Latin Modern Roman}
% ============================================================================
% <<< FIN BLOC CHARTE
% ============================================================================

\hypersetup{
  linkcolor=ScriptoriumAccent,
  citecolor=ScriptoriumEncre,
  urlcolor=ScriptoriumAccent,
  pdftitle={Rapport},
}

% Legendes : numero en gras a la couleur d'encre de la charte, separateur point,
% corps reduit. Place apres le bloc CHARTE, qui definit les couleurs employees.
\captionsetup{font=small, labelfont={bf,color=ScriptoriumEncre}, labelsep=period}
\captionsetup[sub]{font=small, labelfont=bf}

% Sans babel ni polyglossia, les titres des listes restent en anglais
% (Contents, List of Figures, List of Tables). Pour un rapport en francais,
% decommenter les trois lignes suivantes.
% \renewcommand{\contentsname}{Sommaire}
% \renewcommand{\listfigurename}{Table des figures}
% \renewcommand{\listtablename}{Liste des tableaux}

% --- Titres ---
\titleformat{\section}{\normalfont\Large\bfseries\color{ScriptoriumEncre}}{\thesection}{1em}{}
\titleformat{\subsection}{\normalfont\large\bfseries\color{ScriptoriumEncre}}{\thesubsection}{1em}{}
\titleformat{\subsubsection}{\normalfont\normalsize\bfseries\color{ScriptoriumTrait}}{\thesubsubsection}{1em}{}

% --- En-tetes et pieds de page ---
\pagestyle{fancy}
\fancyhf{}
\fancyhead[L]{\small\textit{\leftmark}}
\fancyhead[R]{\small\textit{Rapport}}
\fancyfoot[C]{\thepage}
\renewcommand{\headrulewidth}{0.6pt}
\renewcommand{\headrule}{\hbox to\headwidth{\color{ScriptoriumAccent}\leaders\hrule height\headrulewidth\hfill}}

% ============================================================================
% Encadres semantiques (tcolorbox : l'equivalent le plus portable pour ce type d'encadre,
% disponible dans toute distribution TeX Live standard)
% ============================================================================
\newtcolorbox{resultat}[1][Resultat]{
  colback=ScriptoriumPalUn, colframe=ScriptoriumAccent,
  fonttitle=\bfseries\color{white}, coltitle=white, colbacktitle=ScriptoriumAccent,
  title=#1, boxrule=0.8pt, arc=2pt, left=8pt, right=8pt, top=6pt, bottom=6pt,
}
\newtcolorbox{methode}[1][Methode]{
  colback=ScriptoriumPalDeux, colframe=ScriptoriumTrait,
  fonttitle=\bfseries\color{white}, coltitle=white, colbacktitle=ScriptoriumTrait,
  title=#1, boxrule=0.8pt, arc=2pt, left=8pt, right=8pt, top=6pt, bottom=6pt,
}
\newtcolorbox{avertissement}[1][Avertissement]{
  colback=ScriptoriumPalTrois, colframe=ScriptoriumEncre,
  fonttitle=\bfseries\color{white}, coltitle=white, colbacktitle=ScriptoriumEncre,
  title=#1, boxrule=1.2pt, arc=2pt, left=8pt, right=8pt, top=6pt, bottom=6pt,
}
\newtcolorbox{limite}[1][Limite]{
  colback=ScriptoriumPalQuatre, colframe=ScriptoriumTrait,
  fonttitle=\bfseries\color{ScriptoriumEncre}, coltitle=ScriptoriumEncre, colbacktitle=ScriptoriumPalQuatre,
  title=#1, boxrule=0.5pt, arc=2pt, left=8pt, right=8pt, top=6pt, bottom=6pt,
}
\newtcolorbox{note}[1][Note]{
  colback=ScriptoriumFond, colframe=ScriptoriumTrait,
  fonttitle=\bfseries\color{ScriptoriumEncre}, coltitle=ScriptoriumEncre, colbacktitle=ScriptoriumPalUn,
  title=#1, boxrule=0.4pt, arc=2pt, left=8pt, right=8pt, top=6pt, bottom=6pt,
}

% ============================================================================
% Macros statistiques
% ============================================================================
% p-valeur : \pvalue{0{,}03}. En mode mathematique, ecrire 0{,}03 (accolades) pour un
% separateur decimal francais correct, sinon LaTeX espace les chiffres autour de la virgule.
\newcommand{\pvalue}[1]{\textit{p}~=~#1}
% Intervalle de confiance a 95 % : \CI{borne basse}{borne haute}.
\newcommand{\CI}[2]{IC 95~\%~[#1, #2]}
% Taille d'effet nommee : \effectsize{d}{0{,}6}.
\newcommand{\effectsize}[2]{#1~=~#2}

% --- Page de titre ---
\newcommand{\pagetitre}[4]{%
  % #1 titre, #2 sous-titre, #3 auteur ou equipe, #4 date
  \begin{titlepage}
  \centering
  \vspace*{2.5cm}
  {\Huge\bfseries\color{ScriptoriumEncre} #1\par}
  \vspace{0.5cm}
  {\Large\color{ScriptoriumTrait} #2\par}
  \vspace{3cm}
  {\large\bfseries #3\par}
  \vspace{0.3cm}
  {\large #4\par}
  \vfill
  \textcolor{ScriptoriumAccent}{\rule{\textwidth}{2pt}}
  \end{titlepage}
}

\begin{document}

\pagetitre{Titre du rapport}{Sous-titre ou perimetre de l'etude}{Auteur ou equipe}{Mois annee}

\tableofcontents
\thispagestyle{empty}
\newpage

% Pages liminaires exigees d'un rapport scientifique ou d'un memoire (voir
% livrer/references/mise-en-forme.md). Chaque liste se remplit seule a la
% deuxieme passe de compilation, depuis les \caption du document. Retirer ces
% deux lignes pour un rapport court sans figure ni tableau : une liste vide
% sous son titre est un defaut de remise.
\listoffigures
\listoftables
\thispagestyle{empty}
\newpage

\section{Introduction}

Contexte, objet du rapport et perimetre. Une phrase de portee generale, suivie de
l'observation precise qui motive le rapport.

\begin{note}
Une remarque secondaire, non essentielle au fil du propos, se place ici plutot que dans
le corps du texte.
\end{note}

\section{Methode}

\begin{methode}
Donnees mobilisees, procedure suivie, limites connues des l'origine de la collecte.
\end{methode}

Le corps du texte developpe la methode en continu. Les notations statistiques suivent des
macros dediees : \pvalue{0{,}01}, \CI{1,2}{3,4}, \effectsize{d}{0{,}6}.

\section{Resultats}

\begin{table}[htbp]
\centering
\caption{Exemple de tableau (a remplacer)}
\label{tab:exemple}
\begin{tabular}{@{}lcc@{}}
\toprule
Variable & n & Moyenne \\
\midrule
Premiere serie & 120 & 42,5 \\
Seconde serie  & 118 & 37,1 \\
\bottomrule
\end{tabular}
\end{table}

% ----------------------------------------------------------------------------
% Figure. Bloc laisse en commentaire : le gabarit compile tel quel, sans exiger
% de fichier image. Decommenter et remplacer courbe-mesures par le nom du
% fichier (PDF, PNG ou JPEG), sans extension de preference : xelatex choisit
% alors le format disponible.
%
% Ordre impose : \includegraphics, puis \caption, puis \label. Un \label pose
% avant sa \caption renvoie au numero de section et non au numero de figure.
% Le placeur [htbp] laisse LaTeX optimiser la page ; [H] (paquet float) impose
% la position exacte, a n'employer que si la lecture le demande.
% La legende est autonome : elle nomme l'objet mesure, l'unite et la source.
%
% \begin{figure}[htbp]
% \centering
% \includegraphics[width=0.8\textwidth]{courbe-mesures}
% \caption{Evolution de la variable mesuree sur la periode observee, en unites SI. Source : releves de l'auteur, 2026.}
% \label{fig:mesures}
% \end{figure}
%
% Renvoi depuis le texte, tilde insecable pour eviter la coupure de ligne :
% La dispersion se lit sur la figure~\ref{fig:mesures}, page~\pageref{fig:mesures}.
%
% Figure composee de deux vues (subcaption), numerotees 1a et 1b :
%
% \begin{figure}[htbp]
% \centering
% \begin{subfigure}[b]{0.48\textwidth}
%   \centering
%   \includegraphics[width=\textwidth]{vue-avant}
%   \caption{Avant intervention.}
%   \label{fig:vues-avant}
% \end{subfigure}
% \hfill
% \begin{subfigure}[b]{0.48\textwidth}
%   \centering
%   \includegraphics[width=\textwidth]{vue-apres}
%   \caption{Apres intervention.}
%   \label{fig:vues-apres}
% \end{subfigure}
% \caption{Etat du dispositif avant et apres intervention. Source : photographies de l'auteur.}
% \label{fig:vues}
% \end{figure}
%
% Renvoi a une sous-figure : voir la figure~\ref{fig:vues-apres}.
% ----------------------------------------------------------------------------

\begin{resultat}[Constat principal]
Resume verifiable du resultat central en une phrase, chiffre a l'appui (tableau~\ref{tab:exemple}).
\end{resultat}

\begin{avertissement}
Point de vigilance qui conditionne la lecture du resultat precedent (biais, echantillon, portee).
\end{avertissement}

\section{Discussion}

Interpretation du resultat, mise en regard avec les travaux existants cites en bibliographie.

\begin{limite}
Limites de methode, de portee ou d'echantillon, enoncees sans les minimiser.
\end{limite}

\section{Conclusion}

Synthese des constats et suite eventuelle donnee au travail.

% ============================================================================
% Bibliographie (BibTeX ; voir scripts/citations.py pour produire un .bib verifie)
% ============================================================================
\bibliographystyle{plain}
\bibliography{references}

\appendix
\section{Annexe}

Materiel complementaire (protocole detaille, tableau etendu, glossaire des sigles).

\end{document}
